%% kapa: exam=1 round=30 year=2019 period=1 form=A subject=economics subject_ko=경제학원론 seq=39 number=79
%% source: Q-Net 감정평가사 기출문제 (cst00309) artlSeq=5205205 file=extracted/1008/5205205/2019년 제30회 감정평가사 1차시험 문제지/2019년 제30회 감정평가사 1차시험 문제지/2019년도 제30회 감정평가사 1차 1교시 A형.hwp
%% conversion: kordoc 4.24.51 -> tools/gen1.py (auto); review=proofread-sonnet; flags=none
\begin{question}{79}
실질화폐수요가 이자율과는 음(-)의 관계이고 실질국민소득과는 양(+)의 관계이다. 화폐시장이 균형일 때, 새로운 균형을 이루기 위한 변수들의 변화에 관한 설명으로 옳지 않은 것은? (단, 화폐시장만 고려하며, 화폐수량설이 성립한다. 명목통화량과 물가수준은 외생변수이다.)
\begin{choices}
  \item 물가수준이 하락하는 경우, 이자율이 변하지 않는다면 화폐유통속도도 변하지 않는다.
  \item 물가수준이 하락하는 경우, 이자율이 변하지 않는다면 실질국민소득은 증가한다.
  \item 실질국민소득이 증가하면, 화폐유통속도는 증가한다.
  \item 명목통화량이 감소하는 경우, 실질국민소득이 변하지 않는다면 화폐유통속도는 증가한다.
  \item 명목통화량이 증가하는 경우, 실질국민소득이 변하지 않는다면 이자율은 하락한다.
\end{choices}
\end{question}
\answerall
\answernote{정답표에 1,2,3,4,5 모두 정답으로 표기(전원 정답)}
